\begin{scanpage}{0131}{114}
\chaptertitle{V}{SECONDARY DIFFRACTION}
In the previous chapters, an approximate solution of diffraction problems was carried out which was based on the representation of the fringing field in the form of the sum of the fields from the uniform and nonuniform parts of the surface current. The first field was found by quadratures, and the second field by approximation; it was assumed that the nonuniform part of the current near the discontinuity (edge) of a surface is the same as on a corresponding wedge.

However, the fields found by such a method are actually the fields from the currents flowing, not only on the flat and curved parts of the body's surface, but also to some extent on the geometric extension of these sections. The error in the expressions for the fringing field which is thus introduced is most significant with a glancing incident wave, when the edge zone occupied by the nonuniform part of the current is noticeably broadened, and also with a glancing radiation, when the direction to the observation point forms a small angle with the given section of the surface. In these cases, the results obtained earlier are in need of substantial corrections. We already talked about this briefly in \S\ 6 and \S\ 12.
\end{scanpage}
